% Generated by scripts/generate_assets.py; do not edit.
\begin{table}[!htbp]
\centering
\normalsize
\caption{Horas de base e fadiga: compensação e equivalente de consumo}\label{tab:mapping}
\setlength{\tabcolsep}{4pt}
\begin{tabular}{llrrrr}
\toprule
Cenário & Função & $A_{\rm req}^{40}$ & $CE^{40}$ & $A_{\rm req}^{36}$ & $CE^{36}$ \\
\midrule
Principal: base 44h & Bilateral & 1,57 & 0,18 & 8,38 & -3,42 \\
Principal: base 44h & Acima do pico & 1,53 & 0,22 & 6,52 & -1,66 \\
PNAD integral & Bilateral & 0,07 & 3,10 & 6,80 & -0,79 \\
PNAD integral & Acima do pico & 0,07 & 3,10 & 5,00 & 0,96 \\
Sem fadiga (base 44h) & Bilateral & 2,40 & -0,67 & 7,42 & -2,52 \\
Sem fadiga (base 44h) & Acima do pico & 2,40 & -0,67 & 7,42 & -2,52 \\
\bottomrule
\end{tabular}
\par\smallskip
\begin{minipage}{0.98\linewidth}\normalsize Valores em \%. A limitação inicial a 44h é uma hipótese adicional de mensuração/cumprimento; altera também o denominador da ponte de remuneração horária, que é recalibrada. A base integral preserva todas as horas observadas. A ausência de fadiga mantém a base limitada a 44h e usa $E_Q=1$, $\kappa=0$. Essas alternativas não são observações distintas de horas contratuais.\end{minipage}
\end{table}
